\section{Interpretation and remaining work}
\label{sec:discussion}
% claims: C19 C27 C30 C38 C40 C41 C42 C43 C44 C45 C46
The strongest corrected result is presently about the relationship between
\emph{evaluation population, score construction, and input dependence}. A model
can achieve a high ROC under one comparison population, perform much less well
on unseen users, and improve when only behavior targets contribute to its
score. Those observations motivate mechanistic investigation; they do not
already settle it.

The interpretability method remains useful: compare adaptation-induced states,
fit a sparse approximation, select candidates on a declared discovery group,
measure behavior-loss interventions against realized perturbation controls,
and report attribution with its location-versus-meaning limitation. Its
methodological value depends on executing these checks and exposing negative
or inconclusive outcomes. Existing sparse coding, causal contrasts, and
feasibility mathematics do not become new mathematics merely through application
to insider threats. The mathematical branch is retained as conditional work,
with no new theorem or novelty claim asserted here.

\paragraph{Comparison-model dependencies.}
The separate classical architecture probe fitted with ten r4.2 official
insiders and one r6.2 insider in its nominally normal training pool. For the
local session autoencoders, an audit of actual saved bundles found contaminated
training membership in 94/120 r4.2 bundles and 8/8 r6.2 bundles. All 120 r4.2
bundles have additional official insiders in their test groups. The 26 r4.2
bundles with clean training membership remain candidates for limited reuse;
they do not clear the old benchmark as a whole. The older sequence-baseline
harness shares the faulty fold construction and saves aggregate metrics,
without the checkpoints and per-example predictions needed for score-only
recovery. Refits must save scaler membership, checkpoints, predictions, and
corrected fold definitions before architecture comparisons return.

\paragraph{Scope outside CERT.}
The CERT clock bug does not resolve separate TWOS or LANL issues. The early
TWOS feature ranking had an independent defect; the corrected bounded held-out
study remains inconclusive. LANL has an independent coupling between sampling
and fold hashes. Neither is treated as a cleared replication here. Personality
novelty associations on only five r6.2 insiders remain exploratory and are not
used to explain the detector as a general law.

\paragraph{Next evidence needed.}
The next experiments must name the exact unsupported claim, corrected inputs,
frozen endpoints, eligibility checks, output artifacts, and measured runtime
before scaling. Priority goes to corrected causal selection and controls only
where the retained paper claim needs them; comparator refits must first repair
their training folds. Repeating every historical job is unnecessary, while
relabeling every old output is insufficient. The 46-family recovery ledger
records this distinction. This edition reports the verified subset and the
remaining gaps; it is not a claim that the complete mechanistic program is
finished or ready for submission.

\section{Conclusion}
Corrected labels preserve the main LM adapters but change several evaluated
populations and invalidate some label-dependent experiment choices. Complete
r4.2 behavior-target scoring improves ranking without changing the model or
removing profile context. The interpretability question therefore remains:
which internal pathways contribute to behavior predictions, and are their
measured effects specific to incident-window activity? The current diagnostic
pilot does not establish that specificity. A defensible account keeps this open
question visible while retaining the verified score-level findings.
